\documentclass[12pt]{amsart} 

\input{./latex-preamble/cm-preamble.tex}

%\graphicspath{ {./diagrams/} }

\DeclareMathOperator{\CaCl}{CaCl}
\DeclareMathOperator{\Pic}{Pic}

\newtheorem{lemma}{Lemma}

\begin{document}

{Hartshorne} \hfill {3-16-2025}

\bigskip\bigskip

\begin{center}

{\sc Chapter II.6, Week 8}

\end{center}

\begin{center}

  {Noah Parker} %\footnote{Collaborated with Kalev Martinson and Frank :).}}
    
\end{center}

\bigskip\bigskip

\begin{enumerate}
  \item[6.1.] {
      Let $X$ be a scheme satisfying $(*)$.
      Then $X\times \PP^n$ also satisfies $(*)$, and $\Cl(X\times \PP^n)\cong \Cl(X)\times \ZZ$.
    }
\end{enumerate}

\begin{proof}
  We note that $X\times \PP^n$ is covered by finitely many $X\times U_i\cong X\times \AA^n$.
  $X\times \AA^n$ satisfies $(*)$ by (6.6) and induction on $n$.
  It follows that $X\times \PP^n$ is noetherian, reduced, and regular in codimension one.
  Each $(X\times U_i)\cap(X\times U_j)=X\times (U_i\cap U_j)$ is nonempty, so $X\times \PP^n$ is irreducible, and hence integral.
  It suffices, then, to show that $X\times\PP^n$ is separated.

  Let $\pr_2:X\times \PP^n\to \PP^n$ denote the projection.
  Note that our cover $\{U_i\}$ has $\pr_2^{-1}(U_i)=X\times U_i$.
  The maps $X\times U_i\to U_i\to \Spec \ZZ$ are separated, so that $X\times U_i\to U_i$ is separated.
  It follows that $X\times \PP^n\to \PP^n$ is separated, so $\PP^n\to \Spec \ZZ$ separated gives us that $X\times \PP^n$ is separated.

  Consider the closed subset $H:=X\times V(x_0)\subset X\times \PP^n$.
  It is easily seen to be irreducible of codimension one, so that (6.5) gives the following exact sequence:
  \[
    \ZZ\to \Cl(X\times \PP^n)\to \Cl(X\times U_0)\to 0.
  \] 
  (6.6) and induction on $n$ once again gives us $\Cl(X\times U_0)\cong \Cl(X)$.
  We begin by showing $\ZZ\to \Cl(X\times \PP^n)$ is injective.
  We will first show that $H\in \CaCl (X\times \PP^n)\cong \Pic(X\times \PP^n)$, then show $\sL(H)\mapsto 1\in \ZZ$ under a series of maps.

  For any open affine $\Spec A\cong V\subset X$, $V\times \PP^n\cong \PP^n_A$, and $H|_{V\times U_i}= (x_0)$ for all $i,j$ with $j\neq 0$.
  It follows that $H\in \CaCl (X\times \PP^n)\cong \Pic(X\times \PP^n)$, and $\sL(H)|_{V\times \PP^n}\cong\O_{\PP^n_{A}}(1)$.
  Now fix a particular $\Spec A\cong V\subset X$, and a point $x\in V$.
  %Now, fix a particular $\Spec A\cong V\subset X$, which is noetherian by noetherianity of $X$.
  %Zorn's lemma supplies us with a maximal ideal of $A$.
  %Let $x\in V$ denote the corresponding closed point.
  We can base extend $\Spec k(x)\to \Spec A$ to a map $\PP^n_{k(x)}\to \PP^n_A$.
  Composing with $\PP^n_A\to X\times \PP^n$ gives us a map $f:\PP^n_{k(x)}\to X\times \PP^n$.

  (Ex. 6.8) gives us a map $f^*:\Pic(X\times \PP^n)\to \Pic(\PP^n_{k(x)})\cong \ZZ$ such that 
  \[
    \sL(H)\mapsto \O_{k(x)}^n\mapsto 1\in \ZZ.
  \]
  Then $m\cdot H\mapsto m\in \ZZ$ under this series of maps, so that in particular $m\cdot H\neq 0$.
  Hence, $\ZZ\to \Cl(X\times \PP^n)$ is injective.
  Note as well that the above defines a map $\Cl(X\times \PP^n)\to \Cl(\PP^n_{k(x)})\cong \ZZ$ such that the following diagram, which we have just shown is exact, commutes.
  \[
    \begin{tikzcd}
      0\ar[r]& \ZZ\ar[r]& \Cl(X\times \PP^n)\ar[d]\ar[r] &\Cl(X)\ar[r] &0 \\
             &&\Cl(\PP_{k(x)}^n)\ar[ul,"\cong"]
    \end{tikzcd}
  \] 
  The splitting lemma then implies that, in fact, $\Cl(X\times \PP^n)\cong \Cl(X)\times \ZZ$.
\end{proof}

% \begin{enumerate}
%   \item[6.2.] {
%       \emph{Varieties in Projective Space.} Let $k$ be an algebraically closed field, and let $X$ be a closed subvariety of $\PP^n_k$ which is nonsingular in codimension one (hence satisfies $(*)$).
%       For any divisor $D=\sum n_i Y_i$ on $X$, we define the \emph{degree} of $D$ to be $\sum n_i\deg Y_i$, where $\deg Y_i$ is the degree of $Y_i$, considered as a projective variety itself.
%       \begin{enumerate}
%         \item Let $V$ be an irreducible hypersurface in $\PP^n$ which does not contain $X$, and let $Y_i$ be irreducible components of $V\cap X$.
%           They all have codimension 1 by (I, Ex. 1.8).
%           For each $i$, let $f_i$ be a local equation for $V$ on some open set $U_i$ of $\PP^n$ for which $Y_i\cap U_i\neq \0$, and let $n_i=v_{Y_i}(\ol f_i)$, where $\ol f_i$ is the restriction of $f_i$ to $U_i\cap X$.
%           Then we define the divisor $V.X$ to be $\sum n_iY_i$.
%           Extend by linearity, and show that this gives a well-defined morphism from the subgroup of $\Div \PP^n$ consisting of divisors, none of whose components contain $X$, to $\Div X$.
%         \item If $D$ is a principal divisor on $\PP^n$, and if $D.X$ is defined as in $(a)$, show that $D.X$ is principal on $X$.
%           Thus, we get a homomorphism $\Cl \PP^n\to \Cl X$.
%       \end{enumerate}
%     }
% \end{enumerate}
% 
% \begin{proof}
%   (a) 
% \end{proof}

% \begin{enumerate}
%   \item[6.4.] {
%       Let $k$ be a field of characteristic $\neq 2$.
%       Let $f\in k[x_1,\ldots, x_n]$ be a \emph{square-free} non-constant polynomial.
%       Let $A=k[x_1,\ldots,x_n,z]/(z^2-f)$.
%       Show that $A$ is an integrally closed ring.
%       [\emph{Hint}: The quotient field $K$ of $A$ is just $k(x_1,\ldots, x_n)[z]/(z^2-f)$.
%       It is a Galois extension of $k(x_1,\ldots, x_n)$ with Galois group $\ZZ/2\ZZ$ generated by $z\mapsto -z$.
%       If $\alpha= g+hz\in K$, where $g,h\in k(x_1,\ldots, x_n)$, then the minimal polynomial of $\alpha$ is $t^2-2gt+(g^2-h^2f)$.
%       Now show that $\alpha$ is integral over $k[x_1,\ldots, x_n]$ if and only if $g,h\in k[x_1,\ldots, x_n]$.
%       Conclude that $A$ is the integral closure of $k[x_1,\ldots, x_n]$ in $K$.
%     }
% \end{enumerate}
% 
% \begin{proof}
%   
% \end{proof}

\begin{enumerate}
  \item[6.5.] {
      \emph{Quadric Hypersurfaces.}
      Let $\char k$, and let $X$ be the affine quadric hypersurface $\Spec k[x_0,\ldots, x_n]/(x_0^2+x_1^2+\cdots + x_r^2)$---cf. (I, Ex. 5.12).
      \begin{enumerate}
        \item Show that $X$ is normal if $r\geqs 2$ (use (Ex. 6.4)).
        \item Show by a suitable linear change of coordinates that the equation of $X$ could be written as $x_0x_1=x_2^2 +\cdots + x_r^2$.
          Now, imitate the method of (6.5.2) to show that:
          \begin{enumerate}[label=(\arabic*)]
            \item If $r=2$, then $\Cl X\cong \ZZ/2\ZZ$;
            \item If $r=3$, then $\Cl X\cong \ZZ$ (use (6.6.1) and (Ex. 6.3) above);
            \item If $r\geqs 4$, then $\Cl X=0$.
          \end{enumerate}
        % \item Now let $Q$ be the projective quadric hypersurface in $\PP^n$ defined by the same equation.
        %   Show that:
        %   \begin{enumerate}[label=(\arabic*)]
        %     \item If $r=2$, $\Cl Q\cong \ZZ$, and the class of a hyperplane section $Q.H$ is twice the generator;
        %     \item If $r=3$, $\Cl Q\cong \ZZ\oplus \ZZ$;
        %     \item If $r\geqs 4$, $\Cl Q\cong \ZZ$, generated by $Q.H$.
        %   \end{enumerate}
        % \item Prove Klein's theorem, which says that if $n\geqs r\geqs 4$, and if $Y$ is an irreducible subvariety of codimension 1 on $Q$, then there is an irreducible hypersurface $V\subset \PP^n$ such that $V\cap Q=Y$, with multiplicity one.
        %   In other words, $Y$ is a complete intersection.
          %(\emph{Hint:} First show that)
      \end{enumerate}
    }
\end{enumerate}

\begin{proof}
  (a) Since $r\geqs 1$, $f:=x_1^2+\cdots +x_r^2$ is non-constant.
  We will now show that $f$ is square-free.
  Suppose $f=g^2$ for some $g\in k[x_1,\ldots, x_r]$.
  $f$ is of degree 2, so $g=a+\sum_i a_ix_i$ for some $a,a_i\in k$.
  Then $f=g^2$ implies $a^2=0$, and $2a_ia_j=0$ for all $i\neq j$.
  $\char k\neq 2$ implies $a_ia_j=0$ for all $i\neq j$, so that at most one of the $a_i$ is nonzero.
  Then $f=a_i^2x_i^2$ implies $f$ is a monomial, i.e. $r=1$ and $f=x_1^2$.
  This contradicts $r\geqs 2$, so $f$ is indeed square-free.
  
  Clearly $-f$ is non-constant and square-free as well.
  Thus, we may apply (Ex. 6.4) to show 
  \begin{align*}
    A&= \frac{k[x_0,\ldots,x_r]}{(x_0^2+x_1^2+\cdots +x_r^2)} \\
     &= \frac{k[x_1,\ldots,x_r][x_0]}{(x_0^2-(-f))}
  \end{align*}
  is integrally closed, and hence normal.
  Then $X$ is the spectrum of a normal ring, thus normal itself.

  (b) The map given by $x_0\mapsto \frac i2 (y_0+y_1)$, $x_1\mapsto \frac 12(y_0-y_1)$, and $x_i\mapsto y_i$ for all $1<i\leqs r$ is our desired change of coordinates.
  However, note that it is only well defined if $t^2+1$ splits in $k$, as $x_0$ maps to an element of $k(i)[y_0,y_1]$. 
  
  Clearly $X$ is noetherian and seperated.
  If $r\geqs 2$, then (I,5.12b) gives us $X$ integral and (a) gives us $X$ normal, whence regular in codimension 1.
  If $r=2$, then $X=\Spec k[x,y,z]/(xy-z^2)$, so (1) is exactly (6.5.2).

  Next, let $r=3$.
  The same change of coordinates allows us to write $A=k[x,y,z,w]/(xy-zw)$, so that $X$ is the cone of the surface $Q$ defined in (6.6.1).
  Fix the hyperplane $H:=V(x)\subset \PP^3_k$, which does not contain $Q$.
  (Ex. 6.3b) gives the exact sequence
  \[
    0\to \ZZ\to \Cl Q\to \Cl X\to 0,
  \] 
  where $1\mapsto Q. H\in \Cl Q$.
  $H\cap Q$ corresponds to the hyperplane $\{*\} \times\PP^1$ under the isomorphism $Q\cong \PP^1\times \PP^1$, so $Q.H\mapsto (1,0)$ under $\Cl Q\cong \ZZ\oplus \ZZ$.
  Thus, 
  \[
    \Cl X\cong \frac{\ZZ\oplus \ZZ}{(1,0)\ZZ}\cong \ZZ
  \]
  proves (3).

  Finally, consider the case $r\geqs 4$.
  Use our new coordinates to write $A=k[x_0,\ldots,x_r]/(x_2^2+\cdots+x_r^2-x_0x_1)$, and consider the closed subscheme 
  \[
    Y=\Spec A/(x_0)\cong \Spec k[x_1,x_2,\ldots,x_r]/(x_2^2+\cdots +x_r^2)
  \]
  of $X$.
  $r-2\geqs 2$, so $x_2^2+\cdots+x_r^2$ is irreducible in $k[x_2,\ldots, x_r]$.
  Then 
  \[
    A/(x_0)\cong (k[x_2,\ldots,x_r]/(x_2^2+\cdots+x_r^2))[x_1]
  \]
  is an integral domain, so $x_2^2+\cdots x_r^2$ is irreducible in $k[x_1,\ldots,x_r]$.
  Recalling that $A\cong k[x_0,\ldots,x_r]/(x_0^2+\cdots+x_r^2)$ is quotient by an irreducible ideal as well, Krull's Hauptidealsatz gives us that $\dim A=r$ and $\dim(A/(x_0))=r-1$.
  Since $A/(x_0)$ is integral, $(x_0)\lideal A$ is prime, so (1.8A) gives us $\codim_XY=\htt(x_0)=1$.
  Hence, $Y$ is a prime divisor of $X$.

  By the above, (6.5) supplies us with the exact sequence
  \[
    \ZZ\to \Cl X\to \Cl (X-Y)\to 0,
  \] 
  where $1\mapsto 1\cdot Y$.
  Note that $(x_0)$ is the generic point of $Y$, so the local ring of $Y$ is given by $\O_{X,(x_0)}$.
  The uniformizer is $\O_{X,(x_0)}(x_0)$, so clearly $v_Y(x_0)=1$.
  Clearly $v_{Y'}(x_0)=0$ for all other prime divisors $Y'$, so $Y=(x_0)$ is principle.
  Thus, the exact sequence gives an isomorphism $\Cl X\cong \Cl (X-Y)=\Cl D(x_0)$.
  Recall that $D(x_0)\cong \Spec A_{x_0}$.
  %It suffices, then, to show that $A_{x_0}$ is a UFD.
  We can write
  \begin{align*}
    A_{x_0}&=\frac{k[x_0,x_0^{-1},x_1,x_2,\ldots,x_r]}{x_2^2+\cdots +x_r^2 - x_0x_1} \\
           &\cong \frac{k[x_0,x_0^{-1},x_1,x_2,\ldots,x_r]}{x_0^{-1}(x_2^2+\cdots +x_r^2) - x_0} \\
           &\cong k[x_0,x_0^{-1},x_2,x_3,\ldots,x_r].
  \end{align*}
  This is clearly a UFD, so we have that
  \[
    \Cl X\cong \Cl A_{x_0}=0.
  \] 
  This proves (4).
  % (c) (1) and (3) follow from (b) and (Ex. 6.3), while (2) is exactly (6.6.1).
  % (d) Define $A=k[x_0,\ldots,x_r]/(x_0^2+\cdots +x_r^2)$ as before.
  % $A$ is a noetherian domain with $X=\Spec A$ normal by (a) and $\Cl X=0$ by (b).
  % Hence, $A$ is a UFD by (6.2).
  % Then, for any $n> r$,
  % \begin{align*}
  %   \frac{k[x_0,\ldots,x_n]}{(x_0^2+\cdots +x_r^2)}
  %   &\cong \frac{k[x_0,\ldots,x_{n-1}]}{(x_0^2+\cdots +x_r^2)}[x_n]
  % \end{align*}
  % is a UFD by induction, so that $k[x_0,\ldots,x_n]/(x_0^2+\cdots +x_r^2)$ is a UFD for all $n\geqs r\geqs 4$.

  % Now, fix $Y$, fix $n$, and note that $\Cl \PP^n\cong \ZZ$.
  % Then take some prime divisor $H\subset \PP^n$, so that $\Cl \PP^n = \la H\ra$.
  % (Ex. 5.3) gives us
\end{proof}

% \begin{enumerate}
%   \item[6.10] {
%       \emph{The Grothedieck Group $K(X)$.}
%       Let $X$ be a noetherian scheme.
%       We define $K(X)$ to be the quotient of the free abelian group generated by all coherent sheaves on $X$, by the subgrgoup generated by all expressions $\F-\F'-\F''$, whenever there is an exact sequence $0\to \F'\to\F\to\F''\to 0$ of coherent sheaves on $X$.
%       If $\F$ is a coherent sheaf, we denote by $\gamma(\F)$ its image in $K(X)$.
%       \begin{enumerate}
%         \item If $X=\AA_k^1$, then $K(X)\cong \ZZ$.
%         \item If $X$ is any integral scheme, and $\F$ a coherent sheaf, we define the rank of $\F$ to be $\dim_K\F_\xi$, where $\xi$ is the generic point of $X$, and $K=k(\xi)$ is the function field of $X$.
%           To show that the rank function defines a surjective homomorphism $\rk:K(X)\to \ZZ$.
%         \item If $Y$ is a closed subscheme of $X$, there is an exact sequence
%           \[
%             K(Y)\to K(X)\to K(X-Y)\to 0,
%           \] 
%           where the first map is extension by zero, and the second map is restriction.
%           [\emph{Hint:} For exactness in the middle, show that if $\F$ is a coherent sheaf on $X$ whose support is contained in $Y$, then there is a finite filtration $\F=\F_0\supset \F_1\supset\cdots\supset \F_n=0$, such that each $\F_i/\F_{i+1}$ is an $\O_Y$-module.
%           To show surjectivity on the right, use (Ex. 5.15).]
%       \end{enumerate}
%     }
% \end{enumerate}

\end{document}
